\documentclass[../../main.tex]{subfiles}
\begin{document}

\begin{flushright}
  \footnotesize\textit{【自動文字起こし・要確認】}
\end{flushright}

[解] $P(a, b)$ とし，2接点 $Q(\alpha, \alpha^2), R(\beta, \beta^2)$ とおく．ただし $\alpha < \beta \cdots$ ①とする．

$Q, R$ における接線は，各々
\[
  \begin{cases}
    y = 2\alpha x - \alpha^2 \\
    y = 2\beta x - \beta^2
  \end{cases}
\]
だから，これらの交点が $P$ であり，
\[
  a = \dfrac{\alpha+\beta}{2}, \quad b = \alpha\beta \quad \cdots \text{②}
\]
となる．したがって，
\[
  \begin{aligned}
    \vec{PQ} &=
    \begin{pmatrix} \alpha-a \\ \alpha^2-b
    \end{pmatrix} =
    \begin{pmatrix} \dfrac{\alpha-\beta}{2} \\ \alpha(\alpha-\beta)
    \end{pmatrix} = (\beta-\alpha)
    \begin{pmatrix} -\dfrac{1}{2} \\ -\alpha
    \end{pmatrix} \quad \cdots (*) \\
    \vec{PR} &=
    \begin{pmatrix} \beta-a \\ \beta^2-b
    \end{pmatrix} =
    \begin{pmatrix} \dfrac{-\alpha+\beta}{2} \\ \beta(\beta-\alpha)
    \end{pmatrix} = (\beta-\alpha)
    \begin{pmatrix} \dfrac{1}{2} \\ \beta
    \end{pmatrix}
  \end{aligned}
\]
である．2接線のなす角 $\theta$ として，$\theta = \pi/3$ or $2\pi/3$ だから，
\[
  \tan\theta = \pm \tan\pi/3 = \pm \sqrt{3} \quad \cdots \text{③}
\]
となる．一方，
\[
  \tan\theta = \dfrac{\sin\theta}{\cos\theta} = \dfrac{|\vec{PQ}||\vec{PR}|\sin\theta}{|\vec{PQ}||\vec{PR}|\cos\theta} = \dfrac{2 \Delta PQR}{\vec{PQ}\cdot\vec{PR}} \quad (\because |\vec{PQ}||\vec{PR}| \neq 0)
\]
だから，(*)から
\[
  \tan\theta = \dfrac{|\dfrac{1}{2}\alpha - \dfrac{1}{2}\beta|}{-\dfrac{1}{4} - \alpha\beta} = \dfrac{-\dfrac{1}{2}(\beta-\alpha)}{\dfrac{1}{4} + \alpha\beta} \quad (\because \text{①}) \quad \cdots \text{④}
\]
となる．②から $\alpha, \beta$ が $x$ の2次方程式 $x^2 - 2ax + b = 0$ の2実解であることから，
\[
  \beta-\alpha = 2\sqrt{a^2-b} \quad (a^2-b \geqq 0)
\]
に注意して，②, ④から，
\[
  \pm \sqrt{3} = \dfrac{-\sqrt{a^2-b}}{\dfrac{1}{4}+b}
\]
2乗して良く，
\[
  3\left(b + \dfrac{1}{4}\right)^2 = a^2 - b \iff a^2 - 3\left(b + \dfrac{5}{12}\right)^2 = \dfrac{1}{9} \quad \cdots \text{⑥}
\]
以上から求める軌跡は $a^2 - b \geqq 0 \land \text{⑥}$ だが，①の時 $a^2 - b \geqq 0$ は満たされるので，求めるものは
\[
  x^2 - 3\left(y + \dfrac{5}{12}\right)^2 = \dfrac{1}{9}
\]

\medskip
[解2] 2接線の傾き $\theta_\alpha, \theta_\beta$ とする．これら正負で以下のようになる．
([解1]を流用，$ -\dfrac{\pi}{2} < \theta_\alpha, \theta_\beta < \dfrac{\pi}{2} $)

\begin{figure}[htb]
  \centering
  \begin{tikzpicture}[x=0.95cm, y=0.85cm,
    every node/.style={font=\footnotesize}]
    % Each panel uses actual tangents y=2ux-u^2 at u=alpha,beta.
    % The three choices show positive, opposite, and negative slopes.
    \foreach \aa/\bb/\dx/\title in
    {0.3/1.3/0/{両方正},-0.7/0.7/3.8/{異符号},-1.3/-0.3/7.6/{両方負}} {
      \begin{scope}[shift={(\dx,0)}]
        \pgfmathsetmacro{\px}{(\aa+\bb)/2}
        \pgfmathsetmacro{\py}{\aa*\bb}
        \begin{scope}
          \clip (-1.58,-0.63) rectangle (1.58,2.2);
          \draw[line width=1.1pt, blue!85!black,
          domain=-1.5:1.5, samples=70, smooth]
          plot (\x,{\x*\x});
          \draw[line width=0.9pt, orange!85!black,
          domain=-1.5:1.5] plot (\x,{2*\aa*\x-\aa*\aa});
          \draw[line width=0.9pt, teal!75!black,
          domain=-1.5:1.5] plot (\x,{2*\bb*\x-\bb*\bb});
        \end{scope}
        \draw[->, thin, gray!80!black] (-1.55,0)--(1.58,0);
        \draw[->, thin, gray!80!black] (0,-0.58)--(0,2.18);
        \fill[orange!85!black] (\aa,{\aa*\aa}) circle(1.5pt)
        node[above left=1pt] {$Q$};
        \fill[teal!75!black] (\bb,{\bb*\bb}) circle(1.5pt)
        node[above right=1pt] {$R$};
        \fill[red!80!black] (\px,\py) circle(1.5pt)
        node[below=2pt] {$P$};
        \node[anchor=north] at (0,-0.75) {\title};
      \end{scope}
    }
  \end{tikzpicture}
  \caption{接点 $Q,R$ における接線の傾きの符号による3通りの配置}
  \label{fig:1964_4_tangent_cases}
\end{figure}

いずれの場合にも，これらのなす角 $\theta$ ($0 \leqq \theta \leqq \pi$) は
\[
  \theta = \pi - (\theta_\beta - \theta_\alpha), \quad \theta = \theta_\beta - \theta_\alpha
\]
で与えられ，いずれの場合も
\[
  \tan\theta = \pm \tan(\theta_\beta - \theta_\alpha)
\]
となる．$\theta = \dfrac{\pi}{3}, \dfrac{2\pi}{3}$ だから，
\[
  \pm \sqrt{3} = \pm \tan(\theta_\beta - \theta_\alpha) \quad \cdots \text{P}
\]
ここで $\tan\theta_\alpha = 2\alpha$, $\tan\theta_\beta = 2\beta$ より
\[
  \tan(\theta_\beta - \theta_\alpha) = \dfrac{2(\beta-\alpha)}{1+4\alpha\beta}
\]
だから，①より
\[
  \pm \sqrt{3} = \pm \dfrac{2(\beta-\alpha)}{1+4\alpha\beta} \iff 3 = \dfrac{4(\beta-\alpha)^2}{(1+4\alpha\beta)^2} \quad (\text{Pに合流})
\]
(以下略)

\end{document}
